\chapter{Conclusiones y trabajo futuro}

\section{Conclusiones}
Esta memoria confirma que el uso conjunto de \textit{HDL Coder} y \textit{HDL Verifier} permite establecer un proceso de desarrollo reproducible para pasar de una especificación algorítmica a una implementación en FPGA con evidencia experimental. La validación \textit{FPGA-in-the-Loop} se consolidó como una etapa clave porque habilita pruebas con el mismo banco funcional mientras el procesamiento ocurre en hardware real, reduciendo la incertidumbre entre el comportamiento esperado y el observado en implementación.

Los resultados obtenidos en los casos de estudio muestran que la calidad del diseño final depende menos de la complejidad nominal del algoritmo y más de la disciplina de modelado aplicada desde el inicio. Cuando se definen formatos numéricos coherentes, rutas de datos claramente secuenciadas y límites explícitos de latencia, el flujo de generación se vuelve estable y la síntesis converge con menor retrabajo. En cambio, cuando la especificación numérica queda implícita o se posterga la validación temporal, aparecen desajustes que luego demandan más iteraciones de ajuste.

También se observó que la relación entre recursos y rendimiento no es lineal y debe interpretarse en contexto. Diseños con menor número de operadores pueden presentar latencias elevadas si el patrón de cómputo es secuencial, mientras que arquitecturas más paralelas elevan consumo de recursos, pero reducen la espera por muestra y mejoran la capacidad de operación en tiempo real. Esta tensión entre paralelismo, área y latencia atraviesa todos los ejemplos y refuerza la necesidad de establecer criterios de diseño antes de optimizar casos particulares.

Un aspecto relevante es que la automatización no sustituye el criterio de ingeniería; lo potencia. Las herramientas aceleran tareas repetitivas y mejoran la trazabilidad, pero la decisión sobre qué pipeline usar, qué precisión numérica adoptar y qué compromisos aceptar sigue siendo del diseñador. En ese sentido, la principal contribución práctica de este trabajo no es solo el conjunto de implementaciones obtenidas, sino el marco metodológico para replicarlas y extenderlas de forma controlada.

Finalmente, la evidencia reunida respalda que el enfoque adoptado es útil para docencia, prototipado y desarrollo aplicado. El flujo permite pasar desde ejemplos básicos hasta aplicaciones más exigentes sin cambiar de paradigma de trabajo, conservando continuidad entre modelado, verificación y despliegue. Esta continuidad facilita la formación de criterios técnicos sólidos y reduce la curva de adopción para nuevos proyectos de aceleración en FPGA.

\section{Trabajo futuro}
Como continuidad natural de este trabajo, se propone ampliar el conjunto de algoritmos evaluados con problemas de mayor dimensionalidad y con exigencias temporales estrictas. En particular, resulta pertinente incorporar variantes de control predictivo con horizontes extendidos, redes de inferencia con mayor profundidad y módulos de preprocesamiento orientados a señales ruidosas. Esta expansión permitiría caracterizar con más detalle el punto en que cada estrategia de paralelización deja de ser eficiente para una plataforma determinada.

Otra línea de mejora consiste en fortalecer la automatización del repositorio para transformar el flujo experimental en un entorno de evaluación continua. La integración de scripts de construcción, verificación y recolección de métricas permitiría regenerar tablas y figuras de manera sistemática cada vez que se modifica un modelo. Con ello, la comparación entre versiones pasaría a estar respaldada por evidencia homogénea y no por campañas manuales de prueba.

También se plantea profundizar en métodos de refinamiento numérico orientados a minimizar error sin penalizar área. El análisis de sensibilidad por bloque, la selección dinámica de escalas y la validación cruzada entre precisión y latencia pueden aportar reglas más precisas para decidir formatos \textit{fixed-point}. Este enfoque sería especialmente valioso en algoritmos iterativos, donde pequeños errores de cuantización se acumulan y afectan la convergencia global.

Desde la perspectiva de plataforma, es conveniente evaluar familias de FPGA con diferentes balances entre lógica programable, memoria interna y capacidad de DSP. Una campaña comparativa con criterios comunes de frecuencia, latencia y ocupación permitiría derivar guías de selección de hardware más concretas para proyectos futuros. Asimismo, la exploración de arquitecturas heterogéneas puede abrir nuevas posibilidades cuando conviene distribuir tareas entre procesamiento general y aceleración dedicada.

Por último, se propone extender la documentación metodológica con protocolos de validación por niveles, desde pruebas unitarias de bloques hasta evaluación integrada en lazo cerrado. Este tipo de documentación no solo incrementa la reproducibilidad, sino que facilita la transferencia del flujo a otros equipos de trabajo. Con ello, la memoria deja de ser una fotografía puntual de resultados y se convierte en una base técnica para desarrollos posteriores.
